% Linear functional of local energy. Not unique CHM. Not Einstein.
\documentclass[11pt]{article}
\usepackage[margin=1.05in]{geometry}
\usepackage{amsmath,amssymb,amsthm}
\usepackage[colorlinks=true,linkcolor=blue,citecolor=blue,urlcolor=blue]{hyperref}

\newtheorem*{admission}{Admission}

\title{\textbf{$\delta S$ of Balls Is a Linear Functional\\
of Local Energy. The Kernel Is Not Uniquely CHM.}}
\author{Robert W.\ McGwier\\
\small CTO, Cohere Technology Group\\
\small GitHub: \texttt{n4hy}}
\date{15 August 2026 --- Version 1}

\begin{document}
\maketitle

\begin{abstract}
\noindent
The claim ``entanglement gives linearized gravity'' needs a
linear map from a source to $\delta S$ of balls. This note
builds that map for a free lattice fermion.

One Gaussian potential, $512$ balls of radius $2$ on $N=12$.
$\delta S(\varepsilon)$ versus $\delta S(2\varepsilon)$ has
Pearson $1-7\times 10^{-9}$. $\delta S$ versus the CHM-weighted
energy $\sum (R^2-r^2)\,\delta e$ has $\rho=-0.986$.

A linear $R-r$ weight is slightly better ($R=0.165$ vs $0.168$).
An auditor on $N=10$ still tracks CHM ($\rho=-0.975$) but
\emph{flat} beats CHM. A smooth source does not separate the
kernels.

So: a linear functional, yes. Einstein's kernel, no.
Not $8\pi G$. Not FGHMV in AdS. Not de~Sitter.
\end{abstract}

\section{What this is}

The first-law half of Faulkner--Guica--Hartman--Myers--Van
Raamsdonk, as a lattice measurement, not a bulk Einstein
equation. $n_{\mathrm{occ}}$ was stable. The locked
$\max|\delta S|/\langle S\rangle$ floor failed because $S$ is
the area and $\delta S$ is first order; linearity (C1) is the
real signal check and it passed.

\begin{admission}
I ran the test that can speak to a linear functional. I did
not promote a $0.02$ residual gap into uniqueness, and I did
not write $G_{\mu\nu}=8\pi GT_{\mu\nu}$.
\end{admission}

\end{document}
